%-----------------------------------------------------------
% PDE
%-----------------------------------------------------------
We consider an EEG phantom composed of two contiguous regions—the inner gel volume and the 3D-printed head shell—together with the closed electrical circuit used for signal generation.

\subsubsection{Conductivity.}
The gel is a conductive colloid (saline gelatin) with conductivity on the order of 
$\sigma_{\text{gel}} \approx 0.3{-}0.4~\text{S/m}$ 
\nota{[Open Access Article --- 
A Cost-Effective 3D-Printed Conductive Phantom for EEG Sensing System Validation: Development, Performance Evaluation, and Comparison with State-of-the-Art Technologies]}, 
comparable to physiological saline or brain tissue.

The 3D-printed shell is made from a carbon-doped PLA plastic \nota{add exact model/material and datasheet}, which is substantially more conductive (effective $\sigma_{\text{shell}}$ in the range 3--7~S/m, approximated as $\sim 5$~S/m for a layered print) \nota{Add explanation in subsection of the value... based on material datasheet and link the datasheet from fabricant}. 
Notably, the printed shell exhibits anisotropic conductivity due to the layered fabrication process: measurements from \nota{citar aca el mismo paper que antes} report resistivities of $14.4~\Omega\cdot\text{cm}$ in the planar (X--Y) direction and $27.2~\Omega\cdot\text{cm}$ in the print (Z) direction \nota{check data sheet, has fabricant measurments, and i think they were coincident}, corresponding to conductivities of approximately 6.9~S/m and 3.7~S/m, respectively. 
For modeling simplicity we assume an effective isotropic value $\sigma_{\text{shell}} \approx 5$~S/m as a reasonable average, while acknowledging the underlying anisotropy. 
\nota{The simulation accepts this value as a param so we can make a range of simulations for any coherent value in the range mentioned.}

Both regions are treated as linear, homogeneous, and isotropic conductors.
\nota{ 
add some justif for the shel beeing treated as such. - some relevant sources
Electrical anisotropy and its mitigation in conductive polymers printed by vat photopolymerization - \url{https://www.sciencedirect.com/science/article/pii/S2214860424002677} and \url{https://www.mdpi.com/1424-8220/25/16/4974} talk about anisotropy, processing params, and how they affect, etc. Base case: we leave the processing and possible post-procesing up to the person that prints it, say that the anisotropy can be "limited" with some techniques, here are some [c1], [c2] and mayeb add that we use 5 S/m in the known range because bla bla bla ...
}

Inside the gel, a set of embedded electrode patches (labeled $v_1,\dots,v_9$) serve as voltage-driven sources: they are held at prescribed potentials $u = V_i(t)$ 
(Dirichlet boundary conditions on electrode surfaces), each modelled as a cylinder with the dimensions of a Jack 3.5  audio connector (model SP-35501-OM: \nota{12 mm height, 3.5 mm diameter} \nota{\url{https://www.digikey.fr/fr/products/detail/same-sky-formerly-cui-devices/SP-35501-OM/12822335}}). 

\subsubsection{Electrical insulation.} 
There are no prescribed current sources within the volume; all currents arise solely in response to the applied voltages. 
Any current injected by one internal electrode necessarily returns to the generator through another. 

Although the shell material is conductive, it is not connected to any external circuitry. 
The only interfaces with the outside are the EEG measurement electrodes placed on the shell surface, which --- as in any EEG system --- are assumed to have a very high input impedance. 
This makes the approximation of zero net current flow through them theoretically justified.

Consequently, the phantom can be modeled as an electrically insulated volume. 
This is expressed mathematically through homogeneous Neumann boundary conditions,
\[
    \sigma\,\frac{\partial u}{\partial n} = 0,
\]
on the external surface, implying that the normal current density $J_n = \sigma E_n$ vanishes. 
Physically, the phantom behaves as a closed volume conductor: currents injected inside must return inside, with no leakage into the air (which behaves as an excellent insulator at these field strengths).

External EEG electrodes used for recording do not inject current; they simply sample the surface potential. 
Real EEG amplifiers have input impedances on the order of $1$--$10~\text{M}\Omega$ or higher, so the current they draw is negligible 
\nota{\url{https://www.biosemi.com/publications/pdf/Interference_reduction.pdf}}.
% \nota{\url{https://www.biosemi.nl/forum/viewtopic.php?t=486#:~:text=EEG%20currents%20do%20not%20,not%20influence%20measured%20EEG%20voltages}}

Thus, from the phantom’s perspective, a recording electrode is electrically passive: it does not clamp or perturb the surface potential. 
Accordingly, no Dirichlet or Neumann boundary condition is imposed at measurement locations; they are treated purely as observation points for $u$.


% \textbf{{Ohmic and Frequency Response Behavior.}} We assume the conductive gel behaves Ohmically (purely resistive) over the relevant frequency range (DC to ~100 Hz). Empirical studies on tissue phantoms and biological tissues show that at low frequencies (Hz–kHz) the conductivity is essentially frequency-independent, with very small phase lag (capacitive effects) \url{https://www.mdpi.com/1424-8220/25/16/4974#:~:text=match%20at%20L722%20,characteristics%20of%20EEG%20sensing%20systems}. For example, a gelatin or saline phantom can maintain <2\% variation in conductivity from 0.1 Hz to 1 kHz, indicating minimal dispersion in the EEG band. This justifies treating $\sigma$ as a constant (piece-wise constant) and neglecting any dielectric loss or capacitive reactance --- a crucial assumption for deriving a quasi-static model. We will explicitly justify this approximation using Maxwell’s equations.

\subsubsection{Ohmic and Frequency Response Behavior.}

Over the frequency range of interest (DC to $\sim$100~Hz), the gel medium behaves as a linear, Ohmic conductor with negligible reactive (capacitive or inductive) effects. This is supported by experimental evidence from tissue-mimicking phantoms, which demonstrate frequency-independent conductivity across the EEG band \nota{\url{https://www.mdpi.com/1424-8220/25/16/4974}}. For example, gelatin or saline phantoms typically exhibit less than 2\% variation in $\sigma$ from 0.1 Hz to 1 kHz, and negligible phase shift, confirming that conduction current dominates.

In physical terms, the displacement current density $\partial \mathbf{D}/\partial t$ remains much smaller than the conduction current $\mathbf{J}$ in this regime. This is formally justified by examining the ratio:

\[
\frac{|\partial \mathbf{D}/\partial t|}{|\mathbf{J}|} = \frac{\omega \varepsilon}{\sigma} \ll 1,
\]

as discussed in detail below. Consequently, dielectric storage and capacitive charging are negligible, and the medium can be considered purely resistive. This assumption is central to the quasi-static approximation, and is further justified in Section~\ref{subsection:maxwell} using Maxwell’s equations.

The shell (carbon-doped PLA) is also modeled as Ohmic over the same range. While conductive polymers may exhibit weak frequency dispersion at higher frequencies, empirical phantom studies using Protopasta-type materials confirm phase delay remaining minimal  within EEG-relevant frequencies \nota{https://www.mdpi.com/1424-8220/25/16/4974}. Therefore, we treat both gel and shell as piecewise constant, frequency-independent conductors.


\subsubsection{Summary of Assumptions}
\nota{So, the assumptions made \textbf{up until now} are... and they are based in the system materials and previously explained bla bla bla
Explain here why we do the "deprecitating of the displacement current" section later. or move it before and add it here... whats the best option? double check
}

\begin{itemize}
    \item The phantom’s materials are linear, homogeneous conductors with
    isotropic (effective) conductivity $\sigma(\mathbf{x})$ (piecewise constant: $\sigma_{\text{gel}}$ vs. $\sigma_{\text{shell}}$).
    
    \item No internal current sources or sinks in the volume (all driving currents come from electrode boundary conditions).

    \item No significant polarization or capacitive storage in the bulk (conductive current dominates over displacement current in the frequency range of interest).

    \item Boundary conditions: Dirichlet ($u=V_i(t)$) on internal electrode surfaces; Neumann ($\sigma,\partial_n u=0$) on the external shell surface; no enforced condition for high-impedance measurement electrodes (they effectively float).

    \item Interfaces between gel and shell are in perfect electrical contact (no contact impedance or gaps), so continuity conditions apply. \nota{this we didnt say before... do we keep it here? do i add a text saying those were the already said pnes and we introudce some new ones: etc etc }

\end{itemize}


With this system defined, we proceed to derive the governing partial differential equation (PDE) for the electric potential $u(\mathbf{x},t)$ in the domain $\Omega = \Omega_{\text{gel}} \cup \Omega_{\text{shell}}$.

\subsubsection{From Maxwell’s Equations to governing PDE}
\label{subsection:maxwell}

We begin with Maxwell’s equations in differential form, appropriate for electromagnetic phenomena in the phantom. In particular, we use:

\paragraph*{Gauss’s law}
\[
\nabla\cdot \mathbf{D} = \rho_f
\]
where   
$\mathbf{D}$ is the electric displacement field (C/m$^{2}$) and 
$\rho_f$ is the free charge density (C/m$^{3}$).

There's no free charge in the conductive bulk (any charge separations are negligible in the voulme) therefore $\rho_f\approx 0 \Rightarrow \nabla\cdot \mathbf{D}\approx 0 $.

\paragraph*{Faraday’s law}
\[
\nabla\times \mathbf{E} = -\displaystyle\frac{\partial \mathbf{B}}{\partial t}
\]

where $\mathbf{E}$ is the electric field (V/m) and  $\mathbf{B}$ is the magnetic flux density (T).  


\paragraph*{Ampère’s circuital law}
\[
\nabla\times \mathbf{H} = \mathbf{J} + \displaystyle\frac{\partial \mathbf{D}}{\partial t}
\]
where $\mathbf{H}$ is the magnetic field intensity (A/m), $\mathbf{J}$ is the conduction current density (A/m$^{2}$) and $\partial \mathbf{D}/\partial t$ is the displacement current density (A/m$^{2}$).  

\paragraph*{Gauss’s law for magnetism}
\[
\nabla\cdot \mathbf{B} = 0
\]
where $\mathbf{B}$ is the magnetic flux density (T).  

We now invoke the quasi-static approximation appropriate for bioelectric systems like EEG
\nota{
\url{https://www.frontiersin.org/journals/computational-neuroscience/articles/10.3389/fncom.2010.00138/full},
\url{https://www.researchgate.net/publication/394534573_Reference-Free_EEG_and_Electric_Fields_A_Multi-Scale_Parallax_View}
}. This approximation is justified when both temporal and spatial variations of the fields are sufficiently slow, so that wave propagation effects can be neglected. In our case, the highest relevant frequency (100 Hz) corresponds to a wavelength larger than 300\space km for both materials (as shown in Appendix~\ref{appendix:wavelength}). Even accounting for the reduced propagation speed in biological media (due to high permittivity), the effective wavelength remains many orders of magnitude larger than the phantom’s size (around ~0.2~m, depending on the look-ed at direction). This places the system deep within the electroquasistatic regime, where the electromagnetic fields adjust effectively instantaneously across the volume.

Consequently, the spatial phase variation of the electric field is negligible, and potential changes propagate uniformly. Empirical EEG observations support this: “potentials show no measurable time delay between physically separated regions”
\nota{\url{https://www.researchgate.net/publication/394534573_Reference-Free_EEG_and_Electric_Fields_A_Multi-Scale_Parallax_View}}. This allows us to neglect the wave nature of the field and adopt a quasi-static treatment in which the electric field is determined solely by local sources (none in this case) and boundary conditions.

\paragraph{Neglecting the Displacement Current}
In Ampère’s law, the total current includes both the conduction current density $\mathbf{J}$ and the displacement current density $\partial \mathbf{D}/\partial t$.  At low frequencies in conductive media, the displacement current is negligible compared to the conduction current. The criterion for this is:

\[
\frac{|\partial \mathbf{D}/\partial t|}{|\mathbf{J}|} = \frac{\omega \varepsilon}{\sigma} \ll 1
\]

where $\omega = 2\pi f$ is the angular frequency, $\varepsilon$ is the absolute permittivity of the medium, and $\sigma$ its conductivity.

For biological-like media such as a gel phantom made from solidified gelatin with salt in water, this condition holds robustly. Even a small amount of salt raises conductivity to $\sigma \sim 0.3$ S/m, while water-based media at 100 Hz can have effective permittivities on the order of $10^6 \varepsilon_0$
\nota{sources:}
\url{https://scispace.com/pdf/the-dielectric-properties-of-biological-tissues-i-literature-2m4c4enqvy.pdf#:~:text=survey%20scispace,in%20three%20main%20steps} \url{https://www.sci.utah.edu/~wolters/LiteraturZurVorlesung/Literatur/B:_Basics_Modeling/A:_Quasi-Stationarity/B:_PlonseyHeppner_QuasiStat_BullMath_1967.pdf#:~:text=,8}


This yields:

\[
\omega\varepsilon \approx 2\pi \cdot 100 \cdot 10^6 \cdot \varepsilon_0 \approx 0.0056\ \text{S/m} 
\]

Since $\sigma \gg \omega \varepsilon$, the displacement current can be ignored. The resulting approximation is:

\[
\nabla \times \mathbf{H} \approx \mathbf{J}
\]

This quasi-static form of Ampère’s law is standard in low-frequency biophysics and electrophysiology
\url{https://archiwum.pe.org.pl/articles/2016/12/1.pdf#:~:text=%EF%82%B6%20%EF%82%B6%20%EF%83%91%EF%82%B4%20%EF%80%BD%20%EF%80%AD,that%20the%20displacement%20currents%20can}, and reflects that all relevant current arises from conduction. In this regime, tissues and phantoms “have no significant capacitance: they are either purely resistive or the
frequency of activity is sufficiently low that the capacitance can be neglected”
\url{https://sci-hub.arizonastockbroker.com/10.1016/b978-012372560-8/50028-0} \nota{use una cita directa - reformular desp}

Moreover, the general form of the low-frequency approximation is:

\[
\nabla \times \mathbf{H} = \mathbf{J}_0 + \sigma \mathbf{E}
\]

where $\mathbf{J}_0$ is any impressed (source) current and $\sigma \mathbf{E}$ is the conduction current density \url{https://archiwum.pe.org.pl/articles/2016/12/1.pdf#:~:text=%EF%82%B6%20%EF%82%B6%20%EF%83%91%EF%82%B4%20%EF%80%BD%20%EF%80%AD,that%20the%20displacement%20currents%20can}. In our case, $\mathbf{J}_0 = 0$ in passive regions, so the equation simplifies to:

\[
\nabla \times \mathbf{H} = \sigma \mathbf{E}
\]



\paragraph{Neglecting the inductive effects}
We also assume the electromagnetic induction term in Faraday’s law is negligible. That is, $\nabla\times \mathbf{E} \approx -\partial_t \mathbf{B} \approx 0$. Physically, the currents in our phantom (and in the head) are so slowly varying that they do not produce significant time-varying magnetic flux; any induced electric field from changing $\textbf{B}$ is extremely small compared to the electrostatic field. Under this assumption:
\[
\nabla\times \mathbf{E} \;\approx\; \mathbf{0} \label{faraday-qs}
\]
Therefore, we can define $u(\mathbf{x},t)$ such that
\[
E= ~- \nabla u(\mathbf{x},t).
\]


Iin EEG contexts, the zero-curl condition holds to very high accuracy, enabling the use of an electrostatic potential
\nota{\url{https://www.researchgate.net/publication/394534573_Reference-Free_EEG_and_Electric_Fields_A_Multi-Scale_Parallax_View#:}}
\url{https://www.frontiersin.org/journals/computational-neuroscience/articles/10.3389/fncom.2010.00138/full}


At this point, we have reduced Maxwell’s laws to a simplified electroquasistatic form applicable to the phantom:
\begin{itemize}
    \item $\nabla\cdot \mathbf{D} \approx 0$ 
    \item $\nabla\times \mathbf{E} \approx 0 \Rightarrow \mathbf{E}=-\nabla u$.
    \item $\nabla\times \mathbf{H} = \mathbf{J}$ 
    \item  $\nabla\cdot \mathbf{B}=0$ 
\end{itemize}






\paragraph{Ohm’s Law}
In conductive media, the current density $\mathbf{J}$ is related to the electric field by Ohm’s law:
\[
\textbf{J}(\textbf{x},t) = \sigma(\textbf{x})\, \textbf{E}(\textbf{x},t)
\]

Combining Ohm’s law with $\mathbf{E}=-\nabla u$, we get the current density in terms of the potential:

\[
\textbf{J}(\textbf{x},t) = -\,\sigma(\textbf{x})\,\nabla u(\textbf{x},t) \label{eq:JgradU}
\]

\paragraph{Charge Conservation} 
We now enforce charge conservation in differential form. Maxwell’s equations include the continuity equation $\nabla\cdot \mathbf{J} + \partial \rho_f/\partial t = 0$. In steady or slowly varying conduction, no net charge builds up in any region (any small charge separations that do occur serve only to maintain the electric field in steady state, but do not accumulate over time). Thus $\partial_t \rho_f \approx 0$. Consequently:

\[
\nabla \cdot \textbf{J}(\textbf{x},t) = 0, \qquad\text{for all } \mathbf{x}\in \Omega_{\text{gel}}\cup\Omega_{\text{shell}}\label{eq:continuity}
\]

\nota{This expresses that current is solenoidal – what flows into any small volume must also flow out (no sinks or sources of current density except at the electrodes where we drive the current via boundary conditions).}


\paragraph{Final PDE - change name later}
Substituting Eq. \eqref{eq:JgradU} into Eq. \eqref{eq:continuity} yields the governing equation used for the simulations.
\[
-\nabla \cdot \!\big(\sigma(\mathbf{x})\, \nabla u(\mathbf{x},t)\big) \;=\; 0, \qquad \text{for all } \mathbf{x}\in \Omega_{\text{gel}}\cup\Omega_{\text{shell}}. 
\tag{PDE}
\]


\nota{SOURCES PARA CITAR despues}

{Plonsey, R. & D.B. Heppner (1967). “Considerations of quasi-stationarity in electrophysiological systems.” Bull. Math. Biophysics, 29: 657–664. 
sci.utah.edu
researchgate.net

Plonsey, R. & R. Barr (2007). Bioelectricity: A Quantitative Approach. 3rd ed. (esp. Ch.1-3 on volume conductors and quasi-static approximation)
researchgate.net
sciencedirect.com
.

Hämäläinen, M. et al. (1993). “Magnetoencephalography – Theory, instrumentation, and applications.” Rev. Mod. Phys. 65: 413–497. (Derivation of EEG/MEG forward model under quasi-static assumptions)
frontiersin.org
frontiersin.org
.

Geddes, L.A. & L.E. Baker (1967). “The specific resistance of biological materials.” Med. & Biol. Eng., 5: 271–293. (Classic data on tissue conductivities; shows near-ohmic behavior in low kHz range)
sci.utah.edu
.

Gabriel, S., Lau, R., & C. Gabriel (1996). “The dielectric properties of biological tissues: II. Measurements in the frequency range 10 Hz–20 GHz.” Phys. Med. Biol. 41: 2251–2269. (Comprehensive tissue $\sigma$ and $\varepsilon$ data confirming $\sigma \gg \omega\varepsilon$ at <kHz frequencies)
scispace.com
.

Owda, A.Y. & A. Casson (2019). “Investigating gelatine-based head phantoms for electrophysiology.” (Reports on gel phantom conductivity ~0.3 S/m, stability over frequency, etc.)
mdpi.com
mdpi.com
.

Akor, P. et al. (2025). “A cost-effective 3D-printed conductive phantom for EEG sensing system validation.” Sensors 25(16): 4974 
mdpi.com
mdpi.com
. (Details of a similar phantom: conductive PLA anisotropy, gel conductivities, and frequency response).

BioSemi (2010). Technical note on active electrodes and impedance. (Explains why high amplifier impedance means negligible current draw)
biosemi.nl
.

Nunez, P.L. & R. Srinivasan (2006). Electric Fields of the Brain: The Neurophysics of EEG. 2nd ed. (Discusses volume conduction theory and necessity of quasi-static approximation for defining EEG potentials)
researchgate.net
researchgate.net
.

}